\bta{2008}


\begin{enumerate}
	%\renewcommand{\labelenumi}{\arabic{enumi}.}
	% A(\Alph) a(\alph) I(\Roman) i(\roman) 1(\arabic)
	%设定全局标号series=example	%引用全局变量resume=example
	%[topsep=-0.3em,parsep=-0.3em,itemsep=-0.3em,partopsep=-0.3em]
	%可使用leftmargin调整列表环境左边的空白长度 [leftmargin=0em]
	\item
填空题 (本题满分 30 分，每小题 6 分)	
\begin{enumerate}
	%\renewcommand{\labelenumi}{\arabic{enumi}.}
	% A(\Alph) a(\alph) I(\Roman) i(\roman) 1(\arabic)
	%设定全局标号series=example	%引用全局变量resume=example
	%[topsep=-0.3em,parsep=-0.3em,itemsep=-0.3em,partopsep=-0.3em]
	%可使用leftmargin调整列表环境左边的空白长度 [leftmargin=0em]
	\item
 $\lim\limits _{n \rightarrow \infty} \sum\limits_{k=1}^{n} \frac{k}{n^{2}+k}$ \tk{$  \frac{ 1 }{ 2 }  $} 。


\item
设 $y=y(x)$ 是二阶线性常系数非齐次微分方程 $y^{\prime \prime}+2 y^{\prime}+y=e^{3 x}$ 满足条件 $y(0)=y^{\prime}(0)=0$
的解，且 $y=y(x)$ 有连续的二阶导数，则 $\lim\limits _{x \rightarrow 0} \frac{\ln \left(1+x^{2}\right)}{y(x)}=$ \tk{$ 2 $} 。


\item
已知 $f^{\prime}(x) \cdot \int_{0}^{2} f(x) d x=8$, 且 $f(0)=0$, 则 $\int_{0}^{2} f(x) d x=$ \tk{$\pm 4$} 。


\item
设 $\int_{0}^{\pi}\left[f(x)+f^{\prime \prime}(x)\right] \sin x d x=3, f(\pi)=2$, 则 $f(0)=$ \tk{$ 1 $} 。


\item
在过点 $(0,1)$ 的直线 $y=f(x)$ 中，使得积分值 $\int_{0}^{2}\left[x^{2}-(f(x))^{2}\right] d x$ 达到最大的直线 方程为 \tk{$y=-\frac{3}{4} x+1$} 。
	
	
	
	
\end{enumerate}


\banswer{
\begin{enumerate}
	%\renewcommand{\labelenumi}{\arabic{enumi}.}
	\item
	$\frac12$。用夹逼准则。由 $n^{2}+1\leqslant n^{2}+k\leqslant n^{2}+n\ (1\leqslant k\leqslant n)$ 及
	$\sum\limits_{k=1}^{n}k=\frac{n(n+1)}{2}$ 得
	\[
	\frac{n(n+1)}{2\left(n^{2}+n\right)}\leqslant\sum_{k=1}^{n}\frac{k}{n^{2}+k}
	\leqslant\frac{n(n+1)}{2\left(n^{2}+1\right)} .
	\]
	左端恒等于 $\frac12$，右端 $\to\frac12$，故极限为 $\frac12$。

	\item
	$2$。由 $y(0)=y^{\prime}(0)=0$ 及方程得
	$y^{\prime\prime}(0)=e^{0}-2y^{\prime}(0)-y(0)=1$，故 Taylor 展开
	$y(x)=\frac{y^{\prime\prime}(0)}{2}x^{2}+o\left(x^{2}\right)=\frac{x^{2}}{2}+o\left(x^{2}\right)$。
	又 $\ln\left(1+x^{2}\right)\sim x^{2}$，所以极限 $=\dfrac{x^{2}}{x^{2}/2}=2$。

	\item
	$\pm4$。记 $A=\int_0^2f(x)dx$（常数）。由 $f^{\prime}(x)A=8$ 知
	$f^{\prime}(x)=\frac{8}{A}$ 为常数，结合 $f(0)=0$ 得 $f(x)=\frac{8}{A}x$。代回
	\[
	A=\int_0^2\frac{8}{A}x\,dx=\frac{8}{A}\cdot2=\frac{16}{A}
	\quad\Longrightarrow\quad A^{2}=16,\quad A=\pm4 .
	\]

	\item
	$1$。两次分部积分：
	\[
	\int_0^{\pi}f^{\prime\prime}\sin x\,dx
	=\bigl[f^{\prime}\sin x\bigr]_0^{\pi}-\int_0^{\pi}f^{\prime}\cos x\,dx
	=-\bigl[f\cos x\bigr]_0^{\pi}-\int_0^{\pi}f\sin x\,dx .
	\]
	故 $\int_0^{\pi}\left(f+f^{\prime\prime}\right)\sin x\,dx=-\bigl[f\cos x\bigr]_0^{\pi}
	=f(\pi)+f(0)$。由 $f(\pi)+f(0)=3$、$f(\pi)=2$ 得 $f(0)=1$。

	\item
	$y=-\frac34x+1$。过 $(0,1)$ 的直线设为 $f(x)=kx+1$，则
	\[
	I(k)=\int_0^2\left[x^{2}-(kx+1)^{2}\right]dx
	=\frac83\left(1-k^{2}\right)-4k-2 .
	\]
	由 $I^{\prime}(k)=-\frac{16}{3}k-4=0$ 得 $k=-\frac34$；
	又 $I^{\prime\prime}(k)=-\frac{16}{3}<0$，故为最大值点，所求直线为 $y=-\frac34x+1$。
\end{enumerate}
}

\item 
单选题 (本题满分 30 分，每小题 6 分)
\begin{enumerate}
	%\renewcommand{\labelenumi}{\arabic{enumi}.}
	% A(\Alph) a(\alph) I(\Roman) i(\roman) 1(\arabic)
	%设定全局标号series=example	%引用全局变量resume=example
	%[topsep=-0.3em,parsep=-0.3em,itemsep=-0.3em,partopsep=-0.3em]
	%可使用leftmargin调整列表环境左边的空白长度 [leftmargin=0em]
	\item
	二元函数 $f(x, y)=\left\{\begin{array}{ll}\frac{x y}{x^{2}+y^{2}} & (x, y) \neq(0,0) \\ 0 & (x, y)=(0,0)\end{array}\right.$ 在点 $(0,0)$ 处 \xzanswer{B} 
	
\fourchoices
{连续且两个偏导数都存在}
{不连续但两个偏导数都存在}
{连续但至少有一个偏导数不存在}
{不连续且至少有一个偏导数不存在}

\item 
如图, $f(x), g(x)$ 是两个逐段线性的连续函数, 设 $u(x)=f(g(x))$, 则 $u^{\prime}(1)$ 的值为 \xzanswer{A} 
\begin{figure}[h!]
	\centering
	\includesvg[width=0.23\linewidth]{picture/svg/601-006}
\end{figure}

\fourchoices
{$\frac{3}{4}$}
{$-\frac{3}{4}$}
{$-\frac{1}{12}$}
{$\frac{1}{12}$}



\item
方程 $x e^{-x}=\frac{1}{2 e}$ 的实根数为 \xzanswer{C} 


\fourchoices
{$ 0 $}
{$ 1 $}
{$ 2 $}
{$ 3 $}

\item 
与直线 $L_{1}\left\{\begin{array}{l}x=1 \\ y=-2+t \text { 与直线 } L_{2}: \frac{x+1}{1}=\frac{y+1}{2}=\frac{z-1}{1} \text { 都平行，且过原点 } \\ z=1+t\end{array}\right.$
的平面 $\pi$ 方程是 \xzanswer{B} 


\fourchoices
{$x+y+z=0$}
{$x-y+z=0$}
{$x+y-z=0$}
{$z-y-z=0$}

\item 
设 $x^{2}=\sum\limits_{n=0}^{\infty} a_{n} \cos n x \quad(-\pi \leqslant x \leqslant \pi)$, 则傅里叶系数 $a_{2}=$ \xzanswer{C} 


\fourchoices
{$\frac{-2}{\pi}$}
{$\frac{2}{\pi}$}
{$ 1 $}
{$-1$}




	
	
	
\end{enumerate}


\banswer{
\begin{enumerate}
	%\renewcommand{\labelenumi}{\arabic{enumi}.}
	\item
	\textbf{B}。沿坐标轴 $f(x,0)\equiv0$、$f(0,y)\equiv0$，故
	$f_x(0,0)=f_y(0,0)=0$，两个偏导数都存在。
	但沿 $y=kx$ 趋于原点时 $f=\frac{kx^2}{(1+k^2)x^2}=\frac{k}{1+k^2}$ 与 $k$ 有关
	（取 $k=1$ 得 $\frac12$，取 $k=0$ 得 $0$），故重极限不存在，$f$ 在 $(0,0)$ 不连续。

	\item
	\textbf{A}。由链式法则 $u^{\prime}(1)=f^{\prime}\bigl(g(1)\bigr)\cdot g^{\prime}(1)$。
	由图：$g$ 在 $[0,2]$ 上是从 $(0,6)$ 到 $(2,0)$ 的线段，故 $g(1)=3$，$g^{\prime}(1)=-3$；
	$f$ 在 $[2,8]$ 上是从 $(2,4)$ 出发、斜率为 $-\frac14$ 的线段，而 $g(1)=3\in(2,8)$，
	故 $f^{\prime}(3)=-\frac14$。于是
	\[
	u^{\prime}(1)=\left(-\frac14\right)\times(-3)=\frac34 .
	\]

	\item
	\textbf{C}。令 $\varphi(x)=xe^{-x}$，则 $\varphi^{\prime}(x)=(1-x)e^{-x}$，
	故 $\varphi$ 在 $(-\infty,1)$ 上递增、在 $(1,+\infty)$ 上递减，最大值
	$\varphi(1)=\frac1e$。又 $x\to-\infty$ 时 $\varphi\to-\infty$，
	$x\to+\infty$ 时 $\varphi\to0^{+}$。
	因 $0<\frac{1}{2e}<\frac1e$，水平线 $y=\frac{1}{2e}$ 与两支单调曲线各交一次，
	故实根恰有 $2$ 个。

	\item
	\textbf{B}。$L_1$ 的方向为 $\bm{s}_1=(0,1,1)$，$L_2$ 的方向为 $\bm{s}_2=(1,2,1)$。
	平面与两者都平行，故法向量
	\[
	\bm{n}=\bm{s}_1\times\bm{s}_2=
	\begin{vmatrix}\bm{i}&\bm{j}&\bm{k}\\0&1&1\\1&2&1\end{vmatrix}=(-1,\ 1,\ -1).
	\]
	又平面过原点，故方程为 $-x+y-z=0$，即 $x-y+z=0$。

	\item
	\textbf{C}。$x^2$ 是偶函数，其在 $[-\pi,\pi]$ 上的余弦展开为
	\[
	x^{2}=\frac{\pi^{2}}{3}+\sum_{n=1}^{\infty}\frac{4(-1)^{n}}{n^{2}}\cos nx ,
	\]
	其中 $n\geqslant1$ 时
	$\frac{2}{\pi}\int_0^{\pi}x^{2}\cos nx\,dx=\frac{2}{\pi}\cdot\frac{2\pi(-1)^{n}}{n^{2}}
	=\frac{4(-1)^{n}}{n^{2}}$。
	题中把常数项也记作 $a_0$，故对 $n\geqslant1$ 有 $a_n=\frac{4(-1)^{n}}{n^{2}}$，于是
	$a_2=\frac{4}{4}=1$。
\end{enumerate}
}



\item 
(本题满分 10 分)
求 $\int_{-1}^{2}\left(|x|+2 x^{2}\right) d x$。
\banswer{
	按 $|x|$ 的定义分段积分：
	\[
	\int_{-1}^{2}|x|\,dx=\int_{-1}^{0}(-x)\,dx+\int_{0}^{2}x\,dx
	=\frac12+2=\frac52 ,
	\]
	\[
	\int_{-1}^{2}2x^{2}\,dx=2\left[\frac{x^{3}}{3}\right]_{-1}^{2}=2\left(\frac83+\frac13\right)=6 .
	\]
	故原积分 $=\frac52+6=\frac{17}{2}$。
}


\item 
(本题满分 10 分)
设曲线 $y=x^{2}$ 和直线 $y=t(0<t<1)$ 分别于 $x=0, x=1$ 所围成的面积之和为 $S(t)$, 试判断 $S(t)$ 是否存在最小值, 若存在, 求出其最小值点。

\banswer{
	直线 $y=t$ 与抛物线 $y=x^{2}$ 在 $x=\sqrt{t}\in(0,1)$ 处相交，故面积之和为
	\begin{align*}
	S(t)&=\int_0^{\sqrt t}\left(t-x^{2}\right)dx+\int_{\sqrt t}^{1}\left(x^{2}-t\right)dx\\
	&=\frac23t^{3/2}+\left(\frac13-t+\frac23t^{3/2}\right)
	=\frac43t^{3/2}-t+\frac13 .
	\end{align*}
	求导，$S^{\prime}(t)=2\sqrt{t}-1$，在 $(0,1)$ 内唯一零点 $t=\frac14$；
	且 $t<\frac14$ 时 $S^{\prime}<0$，$t>\frac14$ 时 $S^{\prime}>0$，故 $t=\frac14$ 是最小值点。
	\[
	S\left(\tfrac14\right)=\frac43\cdot\frac18-\frac14+\frac13=\frac16-\frac14+\frac13=\frac14 .
	\]
	所以 $S(t)$ 存在最小值，最小值点为 $t=\frac14$，且 $S(t)_{\min }=\frac{1}{4}$。
}


\item 
(本题满分 10 分)
设 $f(x)$ 在 $[0,1]$ 上可微, $f(1)=2 \int_{0}^{\frac{1}{2}} x f(x) d x$, 求证:
 存在 $\xi \in(0,1)$, 使得 $\xi f(\xi)+f(\xi)=0$。

\banswer{
本题印刷结论为 $\xi f(\xi)+f(\xi)=0$，即 $(1+\xi)f(\xi)=0$。该结论一般不成立。
例如取 $f(x)=1-\frac{9}{11}x$，它在 $[0,1]$ 上恒正且满足条件
\[
f(1)=\frac{2}{11},\qquad
2\int_{0}^{1/2}x f(x)dx
=2\int_{0}^{1/2}x\Bigl(1-\frac{9}{11}x\Bigr)dx
=2\Bigl(\frac18-\frac{3}{88}\Bigr)=\frac{2}{11}=f(1),
\]
但对任意 $\xi\in(0,1)$ 有 $\xi f(\xi)+f(\xi)=(1+\xi)f(\xi)>0$，不可能为零。
故原题印刷有误，其本意应为 $\xi f'(\xi)+f(\xi)=0$。下按修正后的结论给出证明。

由积分中值定理，存在 $\eta\in[0,\frac12]$ 使得
$\int_0^{1/2}x f(x)dx=\frac12\eta f(\eta)$。
代入条件得 $f(1)=2\cdot\frac12\eta f(\eta)=\eta f(\eta)$。
令 $F(x)=x f(x)$，则 $F(\eta)=\eta f(\eta)=f(1)$，而 $F(1)=f(1)$，故 $F(\eta)=F(1)$。
在 $[\eta,1]\subset[0,1]$ 上对 $F$ 应用罗尔定理，存在 $\xi\in(\eta,1)\subset(0,1)$ 使得
$F'(\xi)=0$，即 $\xi f'(\xi)+f(\xi)=0$。
}


\item 
(本题满分 10 分)
设函数 $u=f\left(\ln \sqrt{x^{2}+y^{2}}\right)$, 满足 $\frac{\partial^{2} u}{\partial x^{2}}+\frac{\partial^{2} u}{\partial y^{2}}=\left(x^{2}+y^{2}\right)^{\frac{3}{2}}$, 试求函数 $f$ 的表达式。


\banswer{
	记 $r=\sqrt{x^{2}+y^{2}}$，$t=\ln r$，则 $u=f(t)$ 只与 $r$ 有关。
	极坐标下径向函数的 Laplace 算子为
	$\Delta u=u_{rr}+\frac1r u_r$，而
	\[
	u_r=\frac{f^{\prime}(t)}{r},\qquad
	u_{rr}=\frac{f^{\prime\prime}(t)-f^{\prime}(t)}{r^{2}} ,
	\]
	故
	\[
	\Delta u=\frac{f^{\prime\prime}-f^{\prime}}{r^{2}}+\frac{f^{\prime}}{r^{2}}
	=\frac{f^{\prime\prime}(t)}{r^{2}} .
	\]
	代入条件 $\Delta u=r^{3}$ 得 $f^{\prime\prime}(t)=r^{5}=e^{5t}$。连续积分两次：
	\[
	f^{\prime}(t)=\frac{e^{5t}}{5}+C_{1},\qquad
	f(t)=\frac{1}{25} e^{5 t}+C_{1} t+C_{2} .
	\]
	即 $f(t)=\frac{1}{25}e^{5t}+C_1t+C_2$（$C_1,C_2$ 为任意常数）。
}


\newpage
\item 
(本题满分 10 分)
计算二重积分 $I=\iint_{D} \sqrt{1-y^{2}} d x d y$, 其中 $D$ 为 $x^{2}+y^{2}=1(y>0)$ 与 $y=|x|$ 围成区域。

\banswer{
	$D$ 是单位圆内、位于两直线 $y=x$ 与 $y=-x$ 上方的区域（即 $\theta\in\left[\frac\pi4,\frac{3\pi}{4}\right]$ 的扇形）。
	按 $y$ 先积 $x$：当 $0\leqslant y\leqslant\frac{1}{\sqrt2}$ 时 $-y\leqslant x\leqslant y$，
	当 $\frac{1}{\sqrt2}\leqslant y\leqslant1$ 时 $-\sqrt{1-y^{2}}\leqslant x\leqslant\sqrt{1-y^{2}}$。故
	\begin{align*}
	I&=\int_0^{1/\sqrt2}\sqrt{1-y^{2}}\cdot2y\,dy+\int_{1/\sqrt2}^{1}\sqrt{1-y^{2}}\cdot2\sqrt{1-y^{2}}\,dy\\
	&=\left[-\frac23\left(1-y^{2}\right)^{3/2}\right]_0^{1/\sqrt2}+2\left[y-\frac{y^{3}}{3}\right]_{1/\sqrt2}^{1}\\
	&=\left(\frac23-\frac{\sqrt2}{6}\right)+\left(\frac43-\frac{5\sqrt2}{6}\right)
	=2-\sqrt{2} .
	\end{align*}
	（数值校验：$2-\sqrt2\approx0.5858$，与数值积分一致。）
}



\item 
(本题满分 10 分)
求微分方程 $y^{\prime}+y=y^{2}(\cos x-\sin x)$ 的通解。

\banswer{
	这是 $n=2$ 的 Bernoulli 方程。两端同除 $y^{2}$（$y\not\equiv0$）：
	\[
	y^{-2}y^{\prime}+y^{-1}=\cos x-\sin x .
	\]
	令 $u=y^{-1}$，则 $u^{\prime}=-y^{-2}y^{\prime}$，方程化为一阶线性方程
	\[
	-u^{\prime}+u=\cos x-\sin x,\qquad\text{即}\qquad u^{\prime}-u=\sin x-\cos x .
	\]
	取积分因子 $e^{-x}$，得 $\left(ue^{-x}\right)^{\prime}=e^{-x}(\sin x-\cos x)$。
	注意到 $\left(-e^{-x}\sin x\right)^{\prime}=e^{-x}(\sin x-\cos x)$，故
	\[
	ue^{-x}=-e^{-x}\sin x+C\quad\Longrightarrow\quad u=-\sin x+Ce^{x} .
	\]
	回代 $u=\frac1y$，得通解
	\[
	\frac{1}{y}=Ce^{x}-\sin x,\qquad\text{即}\qquad y=\frac{1}{Ce^{x}-\sin x} ,
	\]
	另有奇解 $y\equiv0$。

	\par\vspace{1mm}\noindent
	\textbf{注：}{\kaishu
	原答案写作 $u(x)=-\sin x\,e^{x}+C$，有误（把 $Ce^{x}$ 误作 $e^{x}\sin x$），
	且只给出 $u$ 而未回代成 $y$。正确结果为 $\frac1y=Ce^{x}-\sin x$；
	代入原方程验证：由 $w=Ce^{x}-\sin x$、$y=\frac1w$ 得
	$y^{\prime}+y=\frac{-w^{\prime}+w}{w^{2}}=\frac{\cos x-\sin x}{w^{2}}=y^{2}(\cos x-\sin x)$，
	恒等成立（数值校验亦相符）。
	}\par
}



\item 
(本题满分 10 分)
求级数 $\sum\limits_{n=1}^{\infty}(-1)^{n-1} n(n+1) x^{n}$ 在收敛区间 $(-1,1)$ 内的核函数 $S(x)$, 并求数项级数 $\sum\limits_{n=1}^{\infty}(-1)^{n-1} \frac{n(n+1)}{3^{n}}$ 的和。

\banswer{
	先求辅助和。在 $|t|<1$ 内由 $\sum\limits_{n=0}^{\infty}t^{n}=\frac{1}{1-t}$ 逐项求导两次得
	$\sum\limits_{n=1}^{\infty}n(n+1)t^{n-1}=\frac{2}{(1-t)^{3}}$，故
	\[
	\sum_{n=1}^{\infty}n(n+1)t^{n}=\frac{2t}{(1-t)^{3}} .
	\]
	把 $t=-x$ 代入，并注意 $(-1)^{n-1}=-(-1)^{n}$：
	\begin{align*}
	S(x)&=\sum_{n=1}^{\infty}(-1)^{n-1}n(n+1)x^{n}=-\sum_{n=1}^{\infty}n(n+1)(-x)^{n}\\
	&=-\frac{-2x}{(1+x)^{3}}=\frac{2 x}{(1+x)^{3}} ,\qquad x\in(-1,1).
	\end{align*}
	再取 $x=\frac13\in(-1,1)$：
	\[
	\sum_{n=1}^{\infty}(-1)^{n-1} n(n+1)\left(\frac{1}{3}\right)^{n}
	=S\!\left(\frac13\right)=\frac{2/3}{\left(4/3\right)^{3}}=\frac23\cdot\frac{27}{64}=\frac{9}{32} .
	\]
}


\item 
(本题满分 10 分)
已知函数 $f(x)$ 在 $[0,+\infty)$ 上可导, $f(0)=1$, 且满足不等式 $f^{\prime}(x)+f(x)-\frac{1}{1+x} \int_{0}^{x} f(t) d t=$ 0 , 求 $f^{\prime}(x)$ 的函数表达式, 并证明不等式:
\[
e^{-x} \leqslant f(x) \leqslant 1 \quad(x \geqslant 0)
\]


\banswer{
令 $F(x)=\int_0^x f(t)dt$，则 $F'(x)=f(x)$，原方程化为
$(1+x)f'(x)+(1+x)f(x)-F(x)=0$。
两边对 $x$ 求导得 $(1+x)f''(x)+(2+x)f'(x)=0$。
令 $q(x)=f'(x)$，则 $(1+x)q'+(2+x)q=0$。注意到
$\bigl((1+x)e^x q(x)\bigr)'=e^x\bigl[(1+x)q'+(2+x)q\bigr]=0$，
故 $(1+x)e^x q(x)=C$，即 $f'(x)=q(x)=C\frac{e^{-x}}{1+x}$。在原方程中令 $x=0$ 得 $f'(0)+f(0)=0$，故 $f'(0)=-1$，$C=-1$。
于是 $f'(x)=-\frac{e^{-x}}{1+x}$，且 $f(x)=1-\int_0^x\frac{e^{-t}}{1+t}dt$。
由 $f'(x)\le0$ 知 $f(x)\le f(0)=1$。又
\[
\int_0^x\frac{e^{-t}}{1+t}dt\le\int_0^x e^{-t}dt=1-e^{-x},
\]
故 $f(x)\ge1-(1-e^{-x})=e^{-x}$。综上 $e^{-x}\le f(x)\le1$（$x\ge0$）。
}



\item 
(本题满分 10 分)
设函数 $f(x), g(x)$ 具有二阶连续导数，曲线积分
\[ 
\oint_{C}\left[y^{2} f(x)+2 y e^{x}+2 y g(x)\right] d x+2[y g(x)+f(x)] d y=0
 \]
其中 $C$ 为平面上任一简单封闭曲线。求上式中的 $f(x)$ 和 $g(x)$ 使得 $f(0)=g(0)=0$。


\banswer{
	曲线积分沿\emph{任一}简单闭曲线都为零，等价于被积表达式为全微分，即 $\frac{\partial Q}{\partial x}=\frac{\partial P}{\partial y}$。记
	\[
	P=y^{2}f(x)+2ye^{x}+2yg(x),\qquad Q=2yg(x)+2f(x) ,
	\]
	则
	\[
	\frac{\partial P}{\partial y}=2yf(x)+2e^{x}+2g(x),\qquad
	\frac{\partial Q}{\partial x}=2yg^{\prime}(x)+2f^{\prime}(x) .
	\]
	两式对一切 $x,y$ 相等，比较 $y$ 的系数与常数项，得方程组
	\[
	f(x)=g^{\prime}(x),\qquad f^{\prime}(x)=e^{x}+g(x) .
	\]
	消去 $f$ 得 $g^{\prime\prime}-g=e^{x}$。齐次解为 $C_{1}e^{x}+C_{2}e^{-x}$；
	因 $1$ 是特征根，特解取 $Axe^{x}$，代入得 $2Ae^{x}=e^{x}$，$A=\frac12$。故
	\[
	g(x)=C_{1} e^{x}+C_{2} e^{-x}+\frac{1}{2} x e^{x},\qquad
	f(x)=g^{\prime}(x)=C_{1} e^{x}-C_{2} e^{-x}+\frac{1}{2}\left(x e^{x}+e^{x}\right).
	\]
	最后用 $f(0)=g(0)=0$ 定常数：
	\[
	g(0)=C_{1}+C_{2}=0,\qquad f(0)=C_{1}-C_{2}+\frac12=0
	\quad\Longrightarrow\quad C_{1}=-\frac14,\ C_{2}=\frac14 .
	\]
	所以
	\[
	g(x)=\frac{1}{2}xe^{x}+\frac{1}{4}\left(e^{-x}-e^{x}\right),\qquad
	f(x)=\frac{1}{2}xe^{x}+\frac{1}{4}\left(e^{x}-e^{-x}\right).
	\]
	（校验：$f=g^{\prime}$ 且 $f^{\prime}-g=e^{x}$ 均成立，$f(0)=g(0)=0$。）

	\par\vspace{1mm}\noindent
	\textbf{注：}{\kaishu
	原答案只给到含 $C_1,C_2$ 的通解，未按题目要求用 $f(0)=g(0)=0$ 定出常数。
	}\par
}

	
	
\end{enumerate}


